
\subsubsection{6.1 Root Cause Analysis}

The uniform degradation of $\rho_j$ across both factual and creative domains (0.01--0.04 vs inferred range 0.75--0.85) demands systematic diagnosis. We frame this as a competing explanations problem: four hypotheses ($H_1$ through $H_4$) are evaluated against empirical evidence, then synthesized into a hierarchical root cause model.

\#\#\#\# 6.1.1 Competing Explanations Framework

\textbf{Evaluation Criteria}:
\begin{enumerate}
\item \textbf{Convergent Evidence}: Does the explanation align with multiple independent data sources (our experiments, literature, sanity checks)?
\item \textbf{Magnitude Fit}: Can the mechanism explain a 20-$80\times$ shift in $\rho_j$?
\item \textbf{Domain Specificity}: Does it predict uniform degradation (both domains affected) or domain-selective degradation (creative only)?
\item \textbf{Testability}: Can the hypothesis be falsified via ablation or calibration studies?
\end{enumerate}

We rank explanations by likelihood (HIGH/MEDIUM/LOW) and assign them to three tiers: Primary (Tier 1), Contributory (Tier 2), Unlikely (Tier 3).

\#\#\#\# 6.1.2 $H_1$: Out-of-Distribution Generalization Gap (Tier 1 - PRIMARY)

\textbf{Mechanism}: DeBERTa-v3-base, trained on SNLI/MNLI (sentence-pair semantic similarity tasks), does not generalize to factual verification tasks (claim-context consistency checking). The model treats claim-context pairs as ''unrelated statements'' rather than ''factual entailment checks,'' defaulting to the ''neutral'' class.

\textbf{Supporting Evidence}:

\begin{enumerate}
\item \textbf{Our Data (Uniform Degradation)}: $\rho_j$ = 0.0354 (factual) and 0.0103 (creative)---both 20-$80\times$ below inferred range. If creative text confuses NLI (original hypothesis), factual text should work normally. Both fail $\rightarrow$ task-general issue, not domain-specific.
\end{enumerate}

\begin{enumerate}
\item \textbf{Our Data (Neutral-Class Dominance)}: Mean $P(\text{neutral})$ = 0.910 (factual), 0.967 (creative). This mechanistically explains low $\rho_j$: when neutral mass $\approx$ 0.9, $\rho_j = (P(\text{entail}) + P(\text{contradict})) / 1.0 \approx 0.09$.
\end{enumerate}

\begin{enumerate}
\item \textbf{Sanity Check Failure}: On manually selected TruthfulQA correct/incorrect answers, $P(\text{entail}|\text{correct}) = 0.11$ (expected > 0.5), $P(\text{contradict}|\text{incorrect}) = 0.08$ (expected > 0.5). The NLI model fails on known ground-truth examples.
\end{enumerate}

\begin{enumerate}
\item \textbf{Literature Triangulation}:
\begin{itemize}
\item Himal-Badu/Prediction-of-Prediction found NLI features dominate over attention mechanisms ($r < 0.1$ for attention-hallucination correlation), with low overall predictive power---consistent with NLI models not calibrated for hallucination detection.
\item Shaguns26/HallucinoGenAI achieved 95\% recall only after threshold tuning from 50\% to 30\%, indicating uncalibrated outputs requiring task-specific adjustment.
\end{itemize}
\end{enumerate}

\textbf{Theoretical Framing}: SNLI/MNLI training objectives optimize for semantic similarity detection (''Do these sentences describe similar situations?''), not factual verification (''Is this claim consistent with this context?''). The distinction:
\begin{itemize}
\item \textbf{Semantic similarity}: ''A dog plays in the park'' $\leftrightarrow$ ''A puppy runs outside'' $\rightarrow$ ENTAILMENT (same event)
\item \textbf{Factual verification}: Claim: ''The president was born in 1980.'' Context: ''Obama was born in 1975.'' $\rightarrow$ CONTRADICTION (inconsistent facts)
\end{itemize}

When DeBERTa processes claim-context pairs with limited lexical overlap (common in factual verification), it defaults to ''neutral'' because the SNLI/MNLI training distribution contains few long-context factual verification examples.

\textbf{Magnitude Fit}: STRONG. A model trained for one task (similarity) applied to another (verification) can easily produce 20-$80\times$ magnitude shifts when the output class distributions differ fundamentally.

\textbf{Domain Specificity}: PREDICTS UNIFORM DEGRADATION. Matches observed pattern (both domains affected equally).

\textbf{Likelihood}: HIGH (primary explanation)

\textbf{Falsifiability}: Fine-tune DeBERTa on FEVER or HotpotQA (1000--5000 factual verification examples). If $\rho_j$ improves to 0.70--0.85, $H_1$ confirmed.

\#\#\#\# 6.1.3 $H_2$: Claim Decomposition Quality (Tier 2 - CONTRIBUTORY)

\textbf{Mechanism}: NLTK sentence tokenization produces sentences that lack clear entailment relationships (sentences $\neq$ logical claims). NLI receives fragmented propositions that are genuinely ''neutral'' relative to context.

\textbf{Supporting Evidence}: Mean 5.8 claims/sample (factual), 6.2 (creative)---reasonable for sentence tokenization but no validation against logical claim boundaries. No claim validation step implemented. No manual inspection of extracted claims.

\textbf{Magnitude Fit}: MEDIUM. Claim quality issues can degrade $\rho_j$, but explaining a 20-$80\times$ shift requires assuming nearly all claims are malformed---plausible but not confirmed.

\textbf{Likelihood}: MEDIUM (contributory, but secondary to $H_1$)

\textbf{Falsifiability}: Compare NLTK vs LLM-based claim extraction (GPT-3.5/GPT-4). If $\rho_j$ improves by >0.10 with LLM extraction, $H_2$ confirmed as contributory.

\#\#\#\# 6.1.4 $H_3$: Context Pairing Strategy (Tier 2 - CONTRIBUTORY)

\textbf{Mechanism}: Using full text as context (vs claim-local windows) creates premise-hypothesis pairs that are too distant for the NLI model to detect entailment. The model defaults to ''neutral'' for long-distance dependencies.

\textbf{Supporting Evidence}: CCP paper does not specify context windowing strategy. DeBERTa-v3-base max sequence length = 512 tokens. Long contexts may truncate or dilute relevant information.

\textbf{Magnitude Fit}: MEDIUM. Context windowing affects signal strength but is unlikely to cause a 20-$80\times$ magnitude shift alone (more likely a 2-$5\times$ degradation).

\textbf{Likelihood}: MEDIUM (requires ablation study to quantify; plausible but not confirmed)

\textbf{Falsifiability}: Test full-text vs $\pm 1, \pm 2, \pm 3$ sentence windows. If optimal window size achieves $\rho_j > 0.70, H_3$ confirmed as contributory.

\#\#\#\# 6.1.5 $H_4$: Temperature/Calibration Issue (Tier 3 - UNLIKELY)

\textbf{Mechanism}: NLI model outputs are overconfident in ''neutral'' predictions due to uncalibrated logits. Temperature scaling could shift probability mass to entailment/contradiction classes.

\textbf{Magnitude Fit}: WEAK. Calibration typically improves metrics by 5--20\%, not 20-$80\times$. A 20-$80\times$ shift suggests a deeper issue than miscalibration.

\textbf{Likelihood}: LOW (magnitude argument: calibration alone cannot explain 20-$80\times$ shift)

\textbf{Falsifiability}: Learn temperature $T$ on validation set to minimize ECE. If calibrated $\rho_j$ improves by <0.10, $H_4$ refuted.

\#\#\#\# 6.1.6 Root Cause Hierarchy

\textbf{Table 5: Root Cause Summary with Evidence Strength}

\begin{table*}[htbp]
\centering
\resizebox{\dimexpr 2\columnwidth+\columnsep\relax}{!}{%
\begin{tabular}{llllll}
\toprule
Tier & Hypothesis & Mechanism & Evidence Strength & Magnitude Fit & Likelihood \\
\midrule
1 & $H_1$: OOD Generalization Gap & SNLI/MNLI $\neq$ factual verification & Uniform degradation + sanity check + literature & 20-$80\times$ shift explained & HIGH \\
2 & $H_2$: Claim Decomposition & Sentence $\neq$ logical claim & Plausible mechanism, no ablation & 20-$80\times$ requires all claims malformed & MEDIUM \\
2 & $H_3$: Context Pairing & Full-text vs local windows & Plausible mechanism, no ablation & 2-$5\times$ shift expected & MEDIUM \\
3 & $H_4$: Temperature/Calibration & Overconfident neutral & Known neural network issue & 5-20\% shift typical & LOW \\
\bottomrule
\end{tabular}
}
\end{table*}

\textbf{Primary Cause (Tier 1)}: Out-of-distribution generalization gap from SNLI/MNLI to factual verification. This is the necessary condition---without NLI calibration, $\rho_j$ cannot reach inferred range.

\textbf{Contributory Factors (Tier 2)}: Claim decomposition quality and context pairing strategy likely amplify the primary issue but are not sufficient to cause 20-$80\times$ degradation alone.

\textbf{Unlikely (Tier 3)}: Temperature/calibration is a modulator (affects ranking, not magnitude) rather than a root cause.

\textbf{Synthesis}: The root cause hierarchy suggests a multiplicative failure model:

$$\rho_j^{\text{observed}} = \rho_j^{\text{baseline}} \times \text{OOD penalty} \times \text{claim quality penalty} \times \text{context penalty}$$

If OOD penalty $\approx$ 0.1 (90\% neutral class), claim quality penalty $\approx$ 0.5, and context penalty $\approx$ 0.5, then:

$$\rho_j^{\text{observed}} = 0.80 \times 0.1 \times 0.5 \times 0.5 = 0.02$$

This matches the observed 0.01--0.04 range, supporting the hierarchical model.

\subsubsection{6.2 The Reproducibility Gap}

The CCP paper reports +0.05--0.10 ROC-AUC improvements on biography generation tasks but does not provide raw $\rho_j$ distributions, NLI calibration diagnostics, claim decomposition methodology, context pairing strategy, or hyperparameters. This lack of detail is not unique to CCP---it reflects a field-wide pattern in hallucination detection research. Papers optimize for novelty (reporting metric improvements) over reproducibility (documenting how to achieve the baseline).

\textbf{Three Explanations}:

\begin{enumerate}
\item Our implementation is wrong: We misinterpreted the CCP mechanism despite following published equations and code patterns from independent implementations.
\end{enumerate}

\begin{enumerate}
\item CCP paper uses undocumented techniques: The authors may have applied NLI fine-tuning, claim filtering, or threshold tuning that are critical for achieving reported metrics but were not documented.
\end{enumerate}

\begin{enumerate}
\item CCP paper metrics are not directly comparable: The reported ROC-AUC improvements may involve additional components (e.g., combining CCP with other features) not described in the method section.
\end{enumerate}

We cannot determine which explanation is correct without access to the original implementation or correspondence with the authors. This ambiguity is the cost of irreproducibility: future researchers cannot build on the work because the baseline cannot be established.

\subsubsection{6.3 Recommendations for Authors}

We adapt Dodge et al. (2019) reproducibility checklist for hallucination detection research:

\textbf{R1: Report Raw Metric Distributions, Not Just Aggregates}

Include median, mean, std dev, min, max for all primary metrics ($\rho_j$, claim-level scores, token probabilities). Provide violin plots or histograms showing full distribution. Include per-domain breakdowns.

\textbf{R2: Validate NLI Calibration on Known Examples}

Test NLI model on 10--20 manually verified entailment/contradiction examples from your target domain. Expected behavior: $P(\text{entail}|\text{correct}) > 0.5, P(\text{contradict}|\text{incorrect}) > 0.5$. If failed, document whether you fine-tuned the NLI model or adjusted thresholds.

\textbf{R3: Document Claim Decomposition with Inter-Method Agreement}

Specify claim extraction method: NLTK sentence tokenization, Spacy segmentation, LLM-based extraction, dependency parsing. Report inter-method agreement: Compare two methods on 50--100 samples and compute Krippendorff's $\alpha > 0.7$. If $\alpha < 0.7$, document which method you selected and why.

\textbf{R4: Provide Public Code with Baseline Replication Notebooks}

Include full implementation (data loaders, NLI inference, metric computation, visualization), baseline replication notebook reproducing your paper's main result in <100 lines of code, unit tests (5--10 examples with expected outputs), and configuration files documenting all hyperparameters.

\subsubsection{6.4 Limitations of This Work}

\textbf{L1 (CRITICAL): Measurement Validity Failure}

$\rho_j$ values 20-$80\times$ lower than inferred range (0.01--0.04 vs 0.75--0.85) invalidate all hypothesis tests. Root cause: DeBERTa-v3-base NLI assigns approximately 90% mass to ''neutral'' class. Mitigation for future work: Fine-tune NLI on FEVER/HotpotQA.

\textbf{L2 (HIGH): Claim Decomposition Method}

NLTK sentence tokenization may not capture logical claims (sentences $\neq$ propositions). No inter-method agreement analysis. Mitigation for future work: Compare claim extraction methods.

\textbf{L3 (HIGH): No Baseline Replication}

Did not replicate CCP on TruthfulQA factual domain before testing creative domain transfer. Cannot validate inferred $\rho_j$ range (0.75--0.85). Mitigation for future work: Replicate CCP ROC-AUC on original dataset.

\textbf{L4 (MEDIUM): Context Pairing Strategy}

Used full-text context instead of claim-local windows ($\pm 2$ sentences). May contribute to neutral-class dominance if long-distance dependencies exceed NLI model capacity. Mitigation for future work: Ablate context window size.

\textbf{L5 (LOW): Dataset as Domain Proxy}

TruthfulQA and WritingPrompts are proxies for factual/creative domains but may not capture all ontology-specific features. Mitigation for future work: Test on multiple dataset pairs.

\textbf{L6 (LOW): Single Model Architecture}

Only tested DeBERTa-v3-base. Alternative NLI models may show different $\rho_j$ distributions. Mitigation for future work: Test alternative NLI models.

\textbf{L7 (NEW): Cannot Confirm CCP Implementation Correctness}

We cannot confirm our CCP implementation matches the original paper without access to authors' code or correspondence. Our implementation, following published equations, produced $\rho_j$ values 20-$80\times$ lower than inferred expectations. This suggests either (a) CCP requires undocumented techniques, or (b) our implementation differs from the original. Mitigation: Contact CCP authors for implementation details or pivot to methods with public implementations.

\subsubsection{6.5 Future Work}

\textbf{Tier 1 (Immediate)}: NLI calibration fixes (fine-tuning on FEVER, alternative models, temperature scaling). Test alternative NLI models and hallucination detection baselines (AGSER, HAD) to determine whether neutral-class dominance is DeBERTa-specific or task-general. Success criterion: $\rho_j > 0.70$ on factual text.

\textbf{Tier 2 (Contingent on Tier 1)}: Re-test ontology sensitivity hypothesis with validated methodology. If $\Delta\rho_j > 0.15$, hypothesis confirmed. If $\Delta\rho_j < 0.05$, hypothesis refuted.

\textbf{Tier 3 (Novel Directions)}: NLI domain adaptation for creative text (train NLI to distinguish ''creative truth'' = narrative consistency vs hallucination). Build creative-factual paired dataset.

\textbf{Tier 4 (Long-Term)}: Hallucination detection reproducibility study (replicate CCP, AGSER, HAD, SelfCheckGPT on common benchmarks). Benchmark for creativity-preserving hallucination detection.

\subsubsection{6.6 Broader Impact}

\textbf{Positive}: Transparent failure documentation prevents field-wide repetition of costly mistakes. Our reproducibility recommendations (R1--R4), adapted from Dodge et al. (2019), could improve hallucination detection research quality if adopted.

\textbf{Negative}: May discourage researchers from building on CCP paper due to replication uncertainty. Could slow progress if interpreted as ''hallucination detection doesn't work'' rather than ''specific implementation needs methodological fixes.''

\textbf{Mitigation}: Frame this work as constructive critique (improve standards) rather than dismissal (abandon method). NLI-based hallucination detection remains a promising direction---it requires higher methodological rigor, not abandonment.
